\section{Discussion}
\label{sec:discussion}

In this Master Thesis, a complete implementation of ideal MHD in SPH-EXA was presented. The verification on standard MHD tests reproduces the reference results of \cite{WissingShen2020} and \cite{Price2018} accurately in structure and relevant metrics: the Alfvén wave reaches second-order convergence, the magnetized Sedov blast maintains energy conservation within $0.3\%$, the reference solution of the Brio-Wu shocktube is reproduced closely while keeping post-shock oscillations low, and the divergence error stays bounded and converges with increasing resolution across all tests. 

The MHD extension adds a computational overhead of $\sim 60\%$ wall time per iteration. Targeted optimization could likely reduce this further. In strong and weak scaling experiments the implementation matches the hydro baseline. The $1600^3$ turbulence simulation serves as a proof-of-concept that SPMHD in SPH-EXA is ready to tackle real astrophysical problems on Exascale computing systems. 

ARSLR achieved what we hoped for: small but distinct improvements over the `compromise' of a reduced resistivity coefficient on both fronts. In the shocktube test it does not under-dissipate, showing slightly less oscillation than the baseline, and in the noisy Alfvén wave it cleans uncorrelated noise more aggressively. At the same time it preserves the resolved structure of the Alfvén wave amplitude better and reduces resistive heating in all tests. It does so at reasonable computational cost of $4.8\%$ increased time per iteration. 

However, the improvements remain small and are not game-changing for the presented tests. In particular, no test in the suite exposes a case where the limiter fails, SLR under-dissipates, and the modulation of SLRB or SLRB2 would be needed. The parallel-component limiter is only partially motivated MHD and other limiter-designs, inspecting the transverse components separately, might still prove useful. It remains to be seen in which relevant scenarios ARSLR might bring real structural benefits and where it will fail. A well functioning and proven SLR scheme, which reliably protects resolved structure from unwanted dissipation, could potentially also open the door to further increase dissipation (e.g. by raising the coefficient $\alpha_B$ beyond 1) to hit real discontinuities harder without paying the price in other places. 

A recurring theme was that magnetic results are dominated by the choice of AV. Small-scale velocity noise leads to increased magnetic dissipation via the signal speed or enters the induction equation through the velocity gradient tensor. In the advection loop, the AV scheme ended up mattering far more than the AR scheme, with the most aggressive AV scheme leading to the best results. This suggests that smooth-flow MHD applications might benefit from a small constant floor of viscous dissipation being added. 

There is likely no optimal combination of AV and AR schemes that wins every race. A very recent paper \cite{Karapiperis2026} illustrates this: the authors implement SPMHD in \textsc{SWIFT} \cite{Schaller2024}, targeting cosmological simulations, and for that return to a variant of the Tricco \& Price switch combined with the Alfvén speed. They make the same diagnosis as us, finding the velocity-based signal speed to be overly dissipative in unconverged, disordered velocity fields. But the cures are different: they remove particle velocity from the signal speed; we remove the resolved part of $\mathbf{B}_{ab}$.

Until the holy grail is found, software should ideally offer users flexibility. 

%The search for better AR schemes is far from over. A very recent SPMHD implementation in \textsc{SWIFT} \cite{Karapiperis2026} illustrates this. Targeting cosmological simulations, the authors find the velocity-based signal speed of \cite{Price2018} to be overly dissipative in the unconverged, disordered velocity fields of coarsely resolved structures, and return to the Tricco \& Price steepness switch \eqref{eq:tp_switch}, combined with the Alfvén speed and a reduced $\alpha_\mathrm{max}$ following \cite{HopkinsRaives2016}. Their diagnosis is the one made in §\ref{sec:ar} and confirmed by the advection loop in §\ref{sec:MHD_loop}: $v_{\mathrm{sig},B}$ is purely kinematic and any unresolved velocity structure switches it on. The cures are complementary. They remove the velocity from the signal speed; we remove the resolved part of $\mathbf{B}_{ab}$, which is what lets ARSLR preserve the mode of the noisy Alfvén wave while removing the noise, something a per-particle scalar cannot do. Where both the field and the velocity are unresolved, ARSLR dissipates at full strength and their objection applies to it as well; we have not tested that regime. The SLRB variants already carry the Tricco \& Price measure as a modulator, so combining reconstruction with an Alfvén-speed signal velocity is a one-line change in SPH-EXA and the natural next experiment.

%The search for better AR schemes is far from over. A very recent SPMHD implementation in \textsc{SWIFT} \cite{Karapiperis2026} illustrates this: targeting cosmological simulations, the authors find the velocity-based signal speed of \cite{Price2018} to be overly dissipative in unconverged, disordered velocity fields, and return to the Tricco \& Price steepness switch \eqref{eq:tp_switch} combined with the Alfvén speed, following \cite{HopkinsRaives2016}. 

%the authors find the velocity-based signal speed of \cite{Price2018} to be overly dissipative in the unconverged, disordered velocity fields of coarsely resolved cosmological structures, and return to the Tricco \& Price steepness switc

%Their diagnosis is the same as ours in §\ref{sec:ar}: $v_{\mathrm{sig},B}$ is purely kinematic and fires on velocity noise. They remove the velocity from the signal speed; we remove the resolved part of $\mathbf{B}_{ab}$. Which approach is preferable likely depends on the application, and combining them is an obvious next step.

% AV and AR closely linked. Cannot separate.
% (Iwasaki & Inutsuka 2011) Riemann solvers? increased cost. 
% That new paper using switches again. 
% Outlook on better limiter, constructed with separate aggregation
% optimization work? 

% \paragraph{No single best configuration.}
% There is no combination of AV and AR schemes that wins in every test. The practical conclusion is that the code should offer the options and that the user must know the problem being simulated. 

% \paragraph{Outlook.}
%\todo{cite the recent paper returning to switches and position SLR relative to it} 
% Beyond dissipation, the natural extensions are non-ideal MHD terms and a kernel and initial-condition setup that preserves lattices, which would remove the noise floor that currently limits several tests.